\bookchapter{第五十四章：專案一——我的記帳本（本地版）}

\chapterintro{歡迎來到你的第一個“正經專案”。這個記帳本雖然不聯網，但五臟俱全：能記一筆賬、能看流水、能按月統計、能刪改。做完它，你就完整走了一遍“資料 邏輯 介面”的全流程。別急著抄程式碼——先想清楚再動手，這是咱們一貫的原則。}

\booksection{一、需求思考：先問五個問題}

動手之前，用“需求五要素”把想法捋清楚：

\begin{center}
\begin{tabularx}{\linewidth}{|X|X|}
\hline
\textbf{要素} & \textbf{答案} \tabularnewline
\hline
Who（給誰用） & 自己用，家人朋友也行 \tabularnewline
What（做什麼） & 記收支、看流水、按月彙總、按分類統計 \tabularnewline
Why（為什麼） & 練手 + 搞清楚錢花哪兒了 \tabularnewline
How（怎麼做） & Kotlin + Compose + Room + MVVM，純本地無後端 \tabularnewline
Constraint（約束） & 一週內完成；最低支援 Android 8；不上架也先按上架標準寫 \tabularnewline
\hline
\end{tabularx}
\end{center}

\begin{center}\small\color{gray}記帳本需求五要素\end{center}

然後拆解成“每一步都能用”的小步驟：(1) 建 Room 資料層（能存一筆賬）；(2) 列表頁（能看流水）；(3) 錄入頁（能加賬）；(4) 刪除與編輯；(5) 月度彙總卡片。每完成一步，App 都是一個“能用的東西”。

\booksection{二、資料層：讓“賬”有個家}

先定義一筆賬長什麼樣——欄位要夠用但別貪多：

\begin{lstlisting}[language=Kotlin]
@Entity(tableName = "records")
data class Record(
    @PrimaryKey(autoGenerate = true) val id: Long = 0,
    val type: String,        // "income" 收入 | "expense" 支出
    val amount: Double,      // 金額（元）
    val category: String,    // 分類：餐飲、交通、工資……
    val note: String = "",   // 備註（可選）
    val date: String         // "2026-08-15"
)
\end{lstlisting}

Dao 里加兩條“統計用”的查詢——這是記帳本的靈魂：

\begin{lstlisting}[language=Kotlin]
@Dao
interface RecordDao {
    @Query("SELECT * FROM records ORDER BY date DESC, id DESC")
    fun getAll(): Flow<List<Record>>

    @Insert
    suspend fun insert(record: Record)

    @Delete
    suspend fun delete(record: Record)

    @Query("SELECT COALESCE(SUM(amount),0) FROM records WHERE type='income' AND date LIKE :month || '%'")
    suspend fun monthIncome(month: String): Double   // month 形如 "2026-08"

    @Query("SELECT COALESCE(SUM(amount),0) FROM records WHERE type='expense' AND date LIKE :month || '%'")
    suspend fun monthExpense(month: String): Double
}
\end{lstlisting}

\begin{tcbinfo}{小貼士：金額為什麼用 Double？}
嚴謹的金融 App 會用 BigDecimal 或“以分為單位的 Int”來避免浮點誤差（0.1+0.2$\neq$0.3 的老問題）。練手專案用 Double 足夠；如果哪天做真金融產品，記得換成“分”為單位——這是無數前輩踩出來的坑。

\end{tcbinfo}

\booksection{三、MVVM：介面與邏輯“分家”}

直接讓介面運算元據庫？能跑，但介面程式碼會又臭又長。MVVM 架構讓各司其職：View（介面）只管畫，ViewModel（邏輯中樞）管資料，Model（Room）管儲存。

View（Compose 介面）

ViewModel（狀態與邏輯）

Model（Room 資料層）

\begin{center}\small\color{gray}圖 5：MVVM 三層結構——資料單向流動，各層互不越界\end{center}

\begin{lstlisting}[language=Kotlin]
class RecordViewModel(private val dao: RecordDao) : ViewModel() {

    // 把 Flow 轉成介面能觀察的狀態
    val records = dao.getAll()
        .stateIn(viewModelScope, SharingStarted.WhileSubscribed(), emptyList())

    // 使用者操作：呼叫 suspend 函式
    fun addRecord(type: String, amount: Double, category: String, note: String, date: String) {
        viewModelScope.launch {
            dao.insert(Record(type = type, amount = amount,
                category = category, note = note, date = date))
        }
    }

    fun deleteRecord(record: Record) {
        viewModelScope.launch { dao.delete(record) }
    }
}
\end{lstlisting}

介面（Compose）裡這樣用：val vm: RecordViewModel = viewModel()，然後 val list by vm.records.collectAsState()——資料一變，列表自動重新整理，你什麼都不用做。這就是 MVVM + Flow + Compose 三劍合璧的威力。

\booksection{四、介面：列表頁 + 錄入頁}

列表頁用 LazyColumn（懶載入列表）——資料再多也不卡：

\begin{lstlisting}[language=Kotlin]
@Composable
fun RecordListScreen(vm: RecordViewModel) {
    val records by vm.records.collectAsState()

    Column {
        // 頂部彙總卡片（收入-支出=結餘）
        SummaryCard(records)

        // 流水列表
        LazyColumn {
            items(records, key = { it.id }) { record ->
                RecordRow(record) { vm.deleteRecord(record) }
            }
        }

        // 懸浮"記一筆"按鈕  彈錄入對話方塊
        FloatingActionButton(onClick = { showAddDialog = true }) {
            Text("+")
        }
    }
}

@Composable
fun RecordRow(record: Record, onDelete: () -> Unit) {
    Row(Modifier.fillMaxWidth().padding(12.dp)) {
        Text(record.category, fontWeight = FontWeight.Bold)
        Spacer(Modifier.weight(1f))
        Text(if (record.type == "income") "+${record.amount}" else "-${record.amount}",
             color = if (record.type == "income") Color(0xFF27AE60) else Color(0xFFC0392B))
    }
}
\end{lstlisting}

錄入用 AlertDialog（對話方塊）最省事：幾個輸入框 + 收入/支出切換 + 分類下拉，點確定就調 vm.addRecord(...)。

\begin{tcbwarn}{注意：輸入校驗別偷懶}
金額輸入框一定要校驗：空值、負數、“abc”這類非法輸入都要攔住（用 toDoubleOrNull() 判斷）。App 崩掉最常見的原因之一就是“把使用者的輸入當成了乖孩子”。記住：使用者輸入的一切，都要當敵人審查。

\end{tcbwarn}

\booksection{五、編輯與刪除：讓資料“能進能出”}

只會加不會刪的記帳本，用三天就廢了。補兩個操作：

\begin{lstlisting}[language=Kotlin]
// Dao 裡補兩條：
@Query("UPDATE records SET amount=:amount, category=:category, note=:note WHERE id=:id")
suspend fun update(id: Long, amount: Double, category: String, note: String)

@Query("DELETE FROM records WHERE id=:id")
suspend fun deleteById(id: Long)

// 介面上：長按列表項彈"編輯/刪除"選單（Compose 用 DropdownMenu 或 AlertDialog）
// 刪除前最好彈確認框——"誤刪"是記帳類 App 使用者最痛恨的事
\end{lstlisting}

互動細節兩條經驗：刪除必須二次確認（防止手滑），編輯儘量預填原值（改一個字比重新輸入省事）。別小看這兩個細節——產品感就是這些細節堆出來的。

\booksection{六、驗收：三層檢查 + 兩個進階挑戰}

\booksection{八、驗收：三層檢查 + 兩個進階挑戰}

專案做完，按“功能 邊界 架構”三層驗收：能不能記、能不能刪、重開 App 資料還在不在（功能）；金額輸 0 或負數會不會崩、轉屏會不會丟資料（邊界）；介面和資料有沒有混在一起、欄位有沒有硬編碼（架構）。

基礎版驗收透過後，給自己加兩道“進階菜”：

\begin{enumerate}
  \item 月度篩選：頂部加月份切換（比如 2026-08 v），列表只顯示當月流水；
  \item 分類餅圖：用 Compose 的 Canvas 或引入輕量圖表庫，把"錢花哪兒了"畫成餅圖——做完這個，記帳本就有點"產品樣"了。
\end{enumerate}

\begin{tcbpractice}{來動手：本週完成它}
按“資料層 ViewModel 列表頁 錄入頁 彙總卡片”的順序動手，每天一步。卡住了就把報錯原樣丟給搜尋引擎或 AI（還記得第七篇《AI程式設計實戰篇》嗎？描述問題時說清“期望行為、實際行為、復現步驟”）。完成後用它記一週賬——自己做的 App 用起來，成就感會推著你繼續往前走。

\end{tcbpractice}

\begin{tcbsummary}{本章小結}
\begin{itemize}
  \item 需求五要素 拆成"每步能用"的小步驟，再動手
  \item 資料層：Record 實體 + Dao（含月份統計 SQL）+ Database
  \item MVVM 分工：View 只管畫，ViewModel 管邏輯，Model 管儲存
  \item Flow + stateIn + collectAsState：資料一變介面自動重新整理
  \item LazyColumn 懶載入列表；AlertDialog 做錄入
  \item 三層驗收（功能/邊界/架構）+ 進階挑戰（月度篩選、分類餅圖）
\end{itemize}

\end{tcbsummary}
